We developed a variance-corrected single-index method that reuses
full-return generators without counting market variance twice (Fig.~\ref{fig:preservation};
Table~\ref{tab:aggregate}). Centering and rescaling allowed the
generated asset paths to be combined with a shared market path without
fitting a second generator to each asset's regression residuals.
Across $423$ non-market assets, the method recovered calibrated
market loadings with low error and retained heavy-tailed return
distributions. On $416$ complete asset histories from $2025$, it
improved the mean KS pass rate over centered naive composition and
brought the median ratio of synthetic to observed variance close to
one (Table~\ref{tab:oos_composition}). Residual-fit methods remained competitive. With pooled thresholds,
the corrected method and JumpHMM-on-residuals had $99\%$ VaR
exceedance rates of $1.12\%$ and $1.16\%$, respectively, compared
with the nominal $1\%$ rate (Table~\ref{tab:var}). The short holdout
did not establish calibration or a precise tail-performance ranking. The results supported reuse of existing
generators as a practical alternative to fitting separate residual
models, with validation still needed for each intended application.
The jump experiments established that the correction also worked when
the generators included extended episodes in extreme-return states
(Table~\ref{tab:jump_ablation}).
The SPY jump settings improved multi-asset temporal fit in training
but worsened it in the $2025$ holdout, showing that the single-asset
benefit did not transfer automatically. Correcting variance also left
dependence among asset-specific residuals unresolved. Independent
residual draws omit relationships that matter for portfolio risk and
can overstate the return from rebalancing when positive residual
covariance is missing. The composition rule therefore provides one
component of a multi-asset simulator: it controls how generator
variance is combined with the market factor, while portfolio-level
applications require a joint residual model and validation of both
cross-asset dependence and temporal behavior. Per-path centering also
removes terminal residual uncertainty over the chosen horizon. The
present evidence therefore supports the constrained construction for
daily-fit comparisons, while cumulative-risk applications require a
method that retains cumulative residual variation.
